% Date : 22/02/2023
% Version : 1
% Auteur : GBU
% Relu : 
% Relecteur : 

% ------------------------------------------------------------ %
% ----------------------  Fiche ?????  ----------------------- %
% Thème : ???
% Classes : ???
% ------------------------------------------------------------ %



% ======================================================= % 
% ============== Gestion de la compilation ============== %
% ======================================================= % 
\ifdefined\mainIsLoaded\else
\RequirePackage{import}
\subimport{../../}{_preambule_CdE_PC}
\begin{document}
\fi
% ======================================================= % 
% ======================================================= % 






% ============================================================================ %
%                            FICHE D'ENTRAÎNEMENT                              %
% ============================================================================ %
% ======================= Métadonnées sur la fiche =========================== %
\titreFicheEntrainement		{Statique des fluides GBU}
\grandTheme					{Thermodynamique}
\numeroFiche				{THM04}
\uniqueID					{WZhafEInvkq} % une chaîne de caractères (a-z, A-Z) aléatoire de 10 caractères.
\nombreColonnesReponses		{2} % 1, 2, 3, etc.
% ============================================================================ %

%%%%%%%%%%%%%%%%%%%%%%%%%%%%%%%%%%%%%%%%%%%%%%%%%%%%%%%%%%%%%%%%%%%%%%%%%%%%%%%%%%%%%%%%%%%%%%
\debutFicheEntrainement                                                                      %
%%%%%%%%%%%%%%%%%%%%%%%%%%%%%%%%%%%%%%%%%%%%%%%%%%%%%%%%%%%%%%%%%%%%%%%%%%%%%%%%%%%%%%%%%%%%%%




% --------------------------------------------------------------------- %
%                              Prérequis                                %
%                             (facultatif)                              %
% --------------------------------------------------------------------- %
% ================== Métadonnées sur le prérequis ===================== %
\avecPrerequis{Y} % Y ou N
% ===================================================================== %
\begin{prerequis}
	EFSF, Poussée d'Archimède
\end{prerequis}
% --------------------------------------------------------------------- %







% ********************************************************************* %
%                           ENTRAÎNEMENT                                % 
% ********************************************************************* %
% ================ Métadonnées sur l'entraînement ===================== %

\pointInterrogation				{N}
\titreEntrainementFacultatif	{Concept de poussée d'Archimède}
\hauteurLargeurCadreReponse		{8mm}{1.5cm}
\dureeResolutionFacultative		{1} % 1, 2, 3 ou 4
\typeEntrainement				{QCM} % Newton ou QCM ou AN ou vide (exercice "normal")
\nombreColonnesQuestions		{1} % vide, 1, 2, 3, etc.
\avecPlusieursQuestions			{Y} % Y ou N
\initialisationEntrainement
% ===================================================================== %



% =============================================================================== %
\debutEntrainement
% =============================================================================== %

% ^^^^^^^^^^^^^^^^^^^^^^^^^^^^^^^^^^^^^^ %
%               Question                 %
% ^^^^^^^^^^^^^^^^^^^^^^^^^^^^^^^^^^^^^^ %
%=====================
\begin{enonce}
Soit deux verres identiques $A$ et $B$. On remplit le verre $A$
jusqu'à une certaine hauteur $h$. Dans le verre $B$, on met quelques
glaçons et on complète avec de l'eau jusqu'à la même hauteur $h$. 

Les masses $m_{A}$ et $m_{B}$ des deux verres vérifient:
	\begin{listeQCM3Colonnes}
	\item $m_{A}<m_{B}$
	\item $m_{A}=m_{B}$
	\item $m_{A}>m_{B}$
	\end{listeQCM3Colonnes}
\end{enonce}

\reponse{\reponseB{}}
\begin{corrige}
La masse $m_{B}$ peut se décomposer en notant $m_{\text{liq}}$la
masse de la partie liquide et $m_{\text{glaçon}}$ celle des glaçons:
\begin{align*}
m_{B} & =m_{\text{liq}}+m_{\text{glaçon}}\\
 & =\rho_{e}\left(V_{\text{tot}}-V_{\text{im}}\right)+m_{\text{glaçon}}
\end{align*}
 en notant $\rho_{e}$ la masse volumique de l'eau, $V_{\text{tot}}$
le volume totale du verre (égal à celui du verre $A$), $V_{\text{im}}$
le volume immergé des glaçons. 

Par ailleurs, l'équilibre mécanique des glaçons donne d'après le PFD:
\[
m_{\text{glaçon}}=\rho_{e}V_{\text{im}}
\]
Ainsi: $m_{B}=m_{A}$.
\end{corrige}
%=====================
\begin{enonce}
Dans le verre $B$, on remplace maintenant les glaçons par des boules
de polystyrène de même masse que les glaçons mais de densité inférieure.

Par à rapport à la hauteur initiale, le niveau dans ce verre:

	\begin{listeQCM3Colonnes}
	\item augmente.
	\item reste le même.
	\item diminue.
	\end{listeQCM3Colonnes}
\end{enonce}

\reponse{\reponseB{}}
\begin{corrige}
Ayant la même masse que les glaçons, le volume immergé avec les boules
ou avec la glace sera le même. La hauteur sera donc identique.
\end{corrige}
%=====================
\begin{enonce}
On remplace cette fois les glaçons par des boules en fer de masse
identique aux glaçons dans le verre $B$. 

Par à rapport à la hauteur initiale, le niveau dans ce verre:

	\begin{listeQCM3Colonnes}
	\item augmente.
	\item reste le même.
	\item diminue. 
	\end{listeQCM3Colonnes}
\end{enonce}

\reponse{\reponseC{}}
\begin{corrige}
On note $V_{\text{sb1}}$ et $V_{\text{sb2}}$ respectivement les
volumes submergés avec les glaçons et avec les boules de fer. On a
les relations:
\begin{align*}
V_{\text{sb1}} & =V_{\text{liq}}+V_{\text{im}}\\
V_{\text{sb2}} & =V_{\text{liq}}+V_{\text{Fe}}
\end{align*}
 De plus, comme les boules de fer sont de même masse que les glaçons:
$m_{\text{glaçon}}=\rho_{e}V_{\text{im}}=m_{\text{Fe}}=\rho_{\text{Fe}}V_{\text{Fe}}$
en notant $\rho_{\text{Fe}}$ la masse volumique du fer et $V_{\text{Fe}}$
leur volume. Ainsi: $V_{\text{Fe}}=\left(\dfrac{\rho_{e}}{\rho_{\text{Fe}}}\right)V_{\text{im}}$.
Ainsi:
\[
V_{\text{sb2}}=V_{\text{liq}}+\left(\dfrac{\rho_{e}}{\rho_{\text{Fe}}}\right)V_{\text{im}}
\]
avec $\dfrac{\rho_{e}}{\rho_{\text{Fe}}}<1$ comme les boules coulent.
Ainsi: $V_{\text{sb2}}<V_{\text{sb1}}$: le niveau diminue.
\end{corrige}


% =============================================================================== %
\finEntrainement
% =============================================================================== %





% ********************************************************************* %
%                           ENTRAÎNEMENT                                % 
% ********************************************************************* %
% ================ Métadonnées sur l'entraînement ===================== %

\pointInterrogation				{N}
\titreEntrainementFacultatif	{Projection de vecteurs}
\hauteurLargeurCadreReponse		{8mm}{1.5cm}
\dureeResolutionFacultative		{1} % 1, 2, 3 ou 4
\typeEntrainement				{Newton} % Newton ou QCM ou AN ou vide (exercice "normal")
\nombreColonnesQuestions		{3} % vide, 1, 2, 3, etc.
\avecPlusieursQuestions			{Y} % Y ou N
\initialisationEntrainement
% ===================================================================== %

Soit un objet au fond de l'eau. 

Déterminer l'expression du vecteur unitaire normal à la surface de l'objet et 
 orienté dans le sens de la force de pression de l'eau sur l'objet.


\begin{center}
\subimport{_images/}{fiche_THM04_objet_fond_v1.tex}
\end{center}

\begin{corrige}
La force de pression est toujours normale à la surface de l'objet
et orienté vers celui-ci.

$\vec{u_{A}}=\dfrac{1}{\sqrt{2}}\left(\vec{e_{x}}-\vec{e_{y}}\right)$;
$\vec{u_{B}}=-\vec{e_{y}}$;
$\vec{u_{C}}=-\cos\left(\dfrac{\pi}{6}\right)\vec{e_{x}}-\sin\left(\dfrac{\pi}{6}\right)\vec{e_{y}}=-\dfrac{1}{2}\left(\sqrt{3}\vec{e_{x}}+\vec{e_{y}}\right)$.

\end{corrige}

% =============================================================================== %
\debutEntrainement
% =============================================================================== %



\begin{enonce}
En A:
\end{enonce}
\reponse{$\dfrac{1}{\sqrt{2}}\left(\vec{e_{x}}-\vec{e_{y}}\right)$}%
\begin{enonce}
En B:
\end{enonce}
\reponse{$-\vec{e_{y}}$}%
\begin{enonce}
En C:
\end{enonce}
\reponse{$-\dfrac{1}{2}\left(\sqrt{3}\vec{e_{x}}+\vec{e_{y}}\right)$}%


% =============================================================================== %
\finEntrainement
% =============================================================================== %









% ********************************************************************* %
%                           ENTRAÎNEMENT                                % 
% ********************************************************************* %
% ================ Métadonnées sur l'entraînement ===================== %

\pointInterrogation				{N}
\titreEntrainementFacultatif	{La pression dans différents repères}
\hauteurLargeurCadreReponse		{8mm}{1.5cm}
\dureeResolutionFacultative		{1} % 1, 2, 3 ou 4
\typeEntrainement				{Newton} % Newton ou QCM ou AN ou vide (exercice "normal")
\nombreColonnesQuestions		{3} % vide, 1, 2, 3, etc.
\avecPlusieursQuestions			{Y} % Y ou N
\initialisationEntrainement
% ===================================================================== %

Exprimer la pression dans l'eau, supposée incompressible et de pression
$\rho$ avec les différents systèmes de coordonnées choisis. On note
$P_{0}$ la pression de l'air à l'interface eau--air.

\begin{center}
\subimport{_images/}{fiche_THM04_pression_reperes_v1.tex}
\end{center}

\begin{corrige}
L'équation fondamentale de la statique des fluides est: 
\[
\overrightarrow{\text{grad}}P=\rho\vec{g}
\]
On projette cette égalité suivant les différents axes:
\begin{enumerate}
\item Suivant l'axe $\left(O_{1}z_{1}\right)$:
\[
\dfrac{\mathrm{d}P}{\mathrm{d}z_{1}}=\rho g\Rightarrow P(z_{1})=\rho gz_{1}+\text{cste}_{1}
\]
à l'interface air--eau: $P(z_{1}=0)=P_{0}=\text{cste}_{1}$:
\[
P\left(z_{1}\right)=P_{0}+\rho gz_{1}
\]
$P_{0}$ la pression de l'air à l'interface eau--air.
\item Suivant l'axe $\left(O_{2}z_{2}\right)$:
\[
\dfrac{\mathrm{d}P}{\mathrm{d}z_{2}}=-\rho g\Rightarrow P(z_{2})=-\rho gz_{2}+\text{cste}_{2}
\]
à l'interface air--eau: $P(z_{2}=H-h)=P_{0}=-\rho g\left(H-h\right)+\text{cste}_{2}$
donc $\text{cste}_{2}=P_{0}+\rho g\left(H-h\right)$
\[
P\left(z_{2}\right)=P_{0}+\rho g\left(H-h-z_{2}\right)
\]
\item Suivant l'axe $\left(O_{3}z_{3}\right)$:
\begin{align*}
\dfrac{\partial P}{\partial z_{3}} & =-\rho g\sin\alpha\\
P(z_{3}) & =-\rho g\sin\alpha z_{3}+\text{cste}_{3}
\end{align*}
Au fond de l'eau, on a $P(z_{3}=0)=P_{0}+\rho gH=\text{cste}_{3}$:
\[
P(z_{3})=\rho g\left(H-z_{3}\sin\alpha\right)+P_{0}
\]
\end{enumerate}
\end{corrige}

Exprimer:

% =============================================================================== %
\debutEntrainement
% =============================================================================== %

\begin{enonce}
$P_{1}(z_{1})$:
\end{enonce}
\reponse{$P_{0}+\rho gz_{1}$}%
\begin{enonce}
$P_{2}(z_{2})$:
\end{enonce}
\reponse{$P_{0}+\rho g\left(H-h-z_{2}\right)$}%

\begin{enonce}
$P_{3}(z_{3})$:
\end{enonce}
\reponse{$\rho g\left(H-z_{3}\sin\alpha\right)+P_{0}$}%


% =============================================================================== %
\finEntrainement
% =============================================================================== %
















% ********************************************************************* %
%                           ENTRAÎNEMENT                                % 
% ********************************************************************* %
% ================ Métadonnées sur l'entraînement ===================== %

\pointInterrogation				{N}
\titreEntrainementFacultatif	{Intégrales multiples}
\hauteurLargeurCadreReponse		{8mm}{3.5cm}
\dureeResolutionFacultative		{1} % 1, 2, 3 ou 4
\typeEntrainement				{Newton} % Newton ou QCM ou AN ou vide (exercice "normal")
\nombreColonnesQuestions		{1} % vide, 1, 2, 3, etc.
\avecPlusieursQuestions			{Y} % Y ou N
\initialisationEntrainement
% ===================================================================== %





% =============================================================================== %
\debutEntrainement
% =============================================================================== %
Calculer les intégrales multiples suivantes:

\begin{enonce}
$\iint_{D}y\cos\left(xy\right)\mathrm{d}x\mathrm{d}y$ avec $D=\left[0,\pi\right]\times\left[0,\pi\right]$
\end{enonce}
\reponse{$\dfrac{1}{\pi}\left(1-\cos\pi^{2}\right)$}%
\begin{corrige}
\begin{align*}
\int_{0}^{\pi}\left(\int_{0}^{\pi}y\cos\left(xy\right)\mathrm{d}x\right)\mathrm{d}y & =\int_{0}^{\pi}\left(\left[\sin\left(xy\right)\right]_{0}^{\pi}\right)\mathrm{d}y\\
 & =\int_{0}^{\pi}\sin\left(\pi y\right)\mathrm{d}y\\
 & =\left[-\dfrac{1}{\pi}\cos\left(\pi y\right)\right]_{0}^{\pi}\\
 & =\dfrac{1}{\pi}\left(1-\cos\pi^{2}\right)
\end{align*}
\end{corrige}
%------------------------------------------------
\begin{enonce}
$\iint_{D}\left(1+4xye^{x^{2}}e^{y^{2}}\right)\mathrm{d}x\mathrm{d}y$
avec $D=\left[0,1\right]\times\left[0,2\right]$
\end{enonce}
\reponse{$2+\left(e-1\right)\left(e^{4}-1\right)$}%
\begin{corrige}
\begin{align*}
\int_{0}^{2}\left(\int_{0}^{1}1+axye^{x^{2}}e^{y^{2}}\mathrm{d}x\right)\mathrm{d}y & =\int_{0}^{2}\left[x+\dfrac{a}{2}ye^{x^{2}}e^{y^{2}}\right]_{0}^{1}\mathrm{d}y\\
 & =\int_{0}^{2}\left(1+\dfrac{a}{2}\left(e^{1}-1\right)ye^{y^{2}}\right)\mathrm{d}y\\
 & =\left[y+\dfrac{a}{4}\left(e^{1}-1\right)e^{y^{2}}\right]_{0}^{\text{2}}\\
 & =2+\dfrac{a}{4}\left(e^{1}-1\right)\left(e^{4}-1\right)\\
 & =2+\left(e^{1}-1\right)\left(e^{4}-1\right)\text{ avec }a=4
\end{align*}
\end{corrige}
%------------------------------------------------
\begin{enonce}
$\iint_{D}xy^{2}\mathrm{d}x\mathrm{d}y$ avec $D=\left\{ \left(x,y\right),0\leq x\leq1-y,\,0\leq y\leq1\right\} $
\end{enonce}
\reponse{$\dfrac{1}{60}$}%
\begin{corrige}
\begin{align*}
\int_{0}^{1}\left(\int_{0}^{1-y}xy^{2}\mathrm{d}x\right)\mathrm{d}y & =\int_{0}^{1}\left[\dfrac{x^{2}y^{2}}{2}\right]_{0}^{1-y}\mathrm{d}y\\
 & =\int_{0}^{1}\dfrac{1}{2}\left(1-y\right)^{2}y^{2}\mathrm{d}y\\
 & =\dfrac{1}{2}\left(\int_{0}^{1}y^{2}+y^{4}-2y^{3}\mathrm{d}y\right)\\
 & =\dfrac{1}{2}\left[\dfrac{y^{3}}{3}+\dfrac{y^{5}}{5}-\dfrac{2}{4}y^{4}\right]_{0}^{1}\\
 & =\dfrac{1}{60}
\end{align*}
\end{corrige}
%------------------------------------------------
\begin{enonce}
$\iint_{D}\mathrm{d}r\cdot r\mathrm{d}\theta$ avec $D=\left\{ \left(r,\theta\right),0\leq r\leq R,\,0\leq\theta\leq2\pi\right\} $
\end{enonce}
\reponse{$\pi R^{2}$}%
\begin{corrige}
\begin{align*}
\int_{0}^{R}\int_{0}^{2\pi}r\mathrm{d}r\mathrm{d}\theta & =\left(\int_{0}^{R}r\mathrm{d}r\right)\left(\int_{0}^{2\pi}\mathrm{d}\theta\right)\\
 & =\left[\dfrac{r^{2}}{2}\right]_{0}^{R}\times\left[\theta\right]_{0}^{2\pi}\\
 & =\pi R^{2}
\end{align*}
\end{corrige}
%------------------------------------------------
\begin{enonce}
$\iint_{D}R\mathrm{d}\theta\cdot R\sin\theta\mathrm{d}\varphi$
avec $D=\left\{ \left(\theta,\varphi\right),\,0\leq\theta\leq\pi,\,0\leq\varphi\leq2\pi\right\} $
\end{enonce}
\reponse{$4\pi R^{2}$}%
\begin{corrige}
\begin{align*}
\int_{0}^{\pi}\int_{0}^{2\pi}R^{2}\sin\theta\mathrm{d}\theta\mathrm{d}\varphi & =R^{2}\left(\int_{0}^{\pi}\sin\theta\mathrm{d}\theta\right)\left(\int_{0}^{2\pi}\mathrel{d}\varphi\right)\\
 & R^{2}\left[-\cos\theta\right]_{0}^{\pi}\times\left[\varphi\right]_{0}^{2\pi}\\
 & =4\pi R^{2}
\end{align*}
\end{corrige}
%------------------------------------------------
\begin{enonce}
$\iiint_{D}\mathrm{d}r\cdot r\mathrm{d}\theta\cdot r\sin\theta\mathrm{d}\varphi$
avec $D=\left\{ \left(r,\theta,\varphi\right),0\leq r\leq R,\,0\leq\theta\leq\pi,\,0\leq\varphi\leq2\pi\right\} $
\end{enonce}
\reponse{$\dfrac{4}{3}\pi R^{3}$}%
\begin{corrige}
\begin{align*}
\int_{0}^{R}\int_{0}^{\pi}\int_{0}^{2\pi}r^{2}\sin\theta\mathrm{d}r\mathrm{d}\theta\mathrm{d}\varphi & =\left(\int_{0}^{R}r^{2}\mathrm{d}r\right)\left(\int_{0}^{\pi}\sin\theta\mathrm{d}\theta\right)\left(\int_{0}^{2\pi}\mathrm{d}\varphi\right)\\
 & =\left[\dfrac{r^{3}}{3}\right]_{0}^{R}\times\left[-\cos\theta\right]_{0}^{\pi}\times\left[\varphi\right]_{0}^{2\pi}\\
 & =\dfrac{4\pi R^{3}}{3}
\end{align*}
\end{corrige}

% =============================================================================== %
\finEntrainement
% =============================================================================== %







% ---------------------------------------------------------------------------- %
%                       Affichage des réponses mélangées                       %
% ---------------------------------------------------------------------------- %
\afficheReponsesMelangees
% ---------------------------------------------------------------------------- %

%%%%%%%%%%%%%%%%%%%%%%%%%%%%%%%%%%%%%%%%%%%%%%%%%%%%%%%%%%%%%%%%%%%%%%%%%%%%%%%%%%%%%%%%%%%%%%
\finFicheEntrainement                                                                        %
%%%%%%%%%%%%%%%%%%%%%%%%%%%%%%%%%%%%%%%%%%%%%%%%%%%%%%%%%%%%%%%%%%%%%%%%%%%%%%%%%%%%%%%%%%%%%%


% ======================================================= % 
% ============== Gestion de la compilation ============== %
% ======================================================= % 
\ifdefined\mainIsLoaded\else
\printReponsesEtCorriges
\end{document}\fi
% ======================================================= % 
% ======================================================= % 